\documentclass[11pt,a4paper]{article}
\usepackage[utf8]{inputenc}
\usepackage[margin=1.0in]{geometry}
\usepackage{xcolor}
\usepackage{hyperref}
\usepackage{tcolorbox}
\usepackage{listings}
\usepackage{enumitem}
\usepackage{titlesec}

% --- Professional Typography ---
\usepackage{inter}
\usepackage{inconsolata}
\renewcommand{\familydefault}{\sfdefault}

% --- Custom Colors ---
\definecolor{tindblue}{HTML}{0F172A}      % Slate 900 (Main)
\definecolor{tindaccent}{HTML}{2563EB}    % Blue 600 (Accent)
\definecolor{tindbg}{HTML}{F8FAFC}        % Slate 50 (Code BG)
\definecolor{tindborder}{HTML}{E2E8F0}    % Slate 200 (Borders)
\definecolor{alertbg}{HTML}{FEF2F2}       % Red 50 (Warning BG)
\definecolor{alertborder}{HTML}{FCA5A5}   % Red 300 (Warning Border)

% --- Page Style Setup ---
\hypersetup{
  colorlinks=true,
  linkcolor=tindaccent,
  urlcolor=tindaccent
}

\titleformat{\section}
  {\color{tindblue}\normalfont\Large\bfseries}
  {\thesection}{1em}{}
\titleformat{\subsection}
  {\color{tindaccent}\normalfont\large\bfseries}
  {\thesubsection}{1em}{}

% --- Code Block Formatting ---
\lstset{
  backgroundcolor=\color{tindbg},
  basicstyle=\ttfamily\small,
  breaklines=true,
  frame=single,
  rulecolor=\color{tindborder},
  keywordstyle=\color{tindaccent}\bfseries,
  showstringspaces=false
}

% --- Custom Boxes ---
\newtcolorbox{importantnote}{
  colback=alertbg,
  colframe=alertborder,
  boxrule=1pt,
  arc=4pt,
  title={\color{red!80!black}\bfseries IMPORTANT SECURITY NOTE},
  fonttitle=\small\bfseries
}

% --- Document Metadata ---
\title{\color{tindblue}\Huge\bfseries Tind SMTP Setup Guide \\ \Large\textmd{Dedicated Support Gmail \& Credentials Configuration}}
\author{\textbf{Tind Development Team}}
\date{\today}

\begin{document}

\maketitle

\begin{abstract}
This document outlines the step-by-step setup procedure to configure a dedicated transaction/support email for the Tind backend system. It details Google Account registration, security settings (App Passwords), and environment variable configuration to enable transaction-triggered notifications.
\end{abstract}

\hrule
\vspace{1.5em}

\section{Introduction}
To send general application emails (such as order updates, payouts, and verification alerts) dynamically from the Tind backend, the server leverages standard SMTP. A dedicated Google Gmail account acts as our SMTP provider. 

\begin{importantnote}
Google blocks external sign-in attempts using standard passwords. To authorize the server, you \textbf{must} enable 2-Step Verification and generate a 16-character \textbf{App Password} as described in Section \ref{sec:app-password}.
\end{importantnote}

\section{Step 1: Create a Dedicated Support Gmail Account}
For customer branding, register a new support address rather than utilizing personal mailboxes.
\begin{enumerate}
    \item Open the \href{https://accounts.google.com/signup}{Google Account Creation Portal}.
    \item Select the option \textbf{``For my personal use''} (recommended and cost-free).
    \item Choose a professional username, for example:
    \begin{itemize}
        \item \texttt{support.tind@gmail.com}
        \item \texttt{tind.notifications@gmail.com}
        \item \texttt{tind.app.alerts@gmail.com}
    \end{itemize}
    \item Complete the standard setup flow, linking a recovery number and email.
\end{enumerate}

\section{Step 2: Enable 2-Step Verification \& Generate App Password}
\label{sec:app-password}
To grant the backend server access to send emails on your behalf:
\begin{enumerate}
    \item Log into your newly created account at the \href{https://myaccount.google.com/}{Google Account Hub}.
    \item Navigate to the \textbf{Security} tab from the left sidebar.
    \item Locate the section titled \textbf{``How you sign in to Google''}.
    \item Click on \textbf{2-Step Verification} and follow the prompts to activate it.
    \item Once enabled, return to the \textbf{2-Step Verification} settings page, scroll to the bottom, and select \textbf{App passwords}.
    \item Input a descriptor (e.g., \texttt{Tind Backend Server}) and press \textbf{Create}.
    \item Copy the generated \textbf{16-character password} (e.g., \texttt{abcd efgh ijkl mnop}). Save it safely, as it will not be displayed again.
\end{enumerate}

\section{Step 3: Update Backend Configuration}
Open the environment file at \texttt{tind-backend/backend/.env} and insert the credentials (ensuring you remove any spaces from the 16-character App Password):

\begin{lstlisting}[language=bash]
# --- SMTP Email Integration ---
SMTP_HOST=smtp.gmail.com
SMTP_PORT=587
SMTP_USER=tind.notifications@gmail.com
SMTP_PASS=abcdefghijklmnop
SMTP_FROM="Tind Support" <tind.notifications@gmail.com>
\end{lstlisting}

\section{Step 4: Verification and Maintenance}
When the server is booted using \texttt{npm run dev} or \texttt{node server.js}:
\begin{itemize}
    \item In \textbf{Development mode} (\texttt{NODE\_ENV=development}), configuring these variables triggers live outbound emails. In their absence, email contents are outputted directly to the server logs.
    \item In \textbf{Production mode}, ensure port \texttt{587} is open on your host firewall to establish secure TLS SMTP handshakes.
\end{itemize}

\end{document}
